\chapter*{Notation}\label{ch:notation}

In this book we employ notation from a variety of mathematical disciplines,
summarized below. The reader is encouraged to consult this list as a reference
as they read the various chapters of the book. We have tried to keep their usage
as coherent as possible across chapters.

As a first example, for a natural number \(n\) the set \(\{1, 2,
\dots, n\}\) is denoted \([n]\).

\section*{Linear Algebra Notation}
\begin{itemize}
    \item Scalars (either deterministic or random) are non-bold,
        e.g., \(x\), \(X\).
    \item Vectors (either deterministic or random) are lower-case
        bold, e.g., \(\vx\), \(\vpi\). Entries of vectors are denoted
        \(x_{i}\) or \(\pi_{i}\), alternatively \((\vx)_{i}\).
        Generic terms (which can be scalars, vectors, matrices, or
        higher-order tensors) are also lower-case bold.
    \item Matrices (either deterministic or random) are capital bold,
        e.g., \(\vX, \vPi\). Entries of matrices are denoted
        \(X_{ij}\) or \(\Pi_{ij}\), alternatively \((\vX)_{ij}\).
        Columns of matrices are lower-case bold vectors, e.g.,
        \(\vx_{i}\) is the \(i\th\) column of \(\vX\), unless otherwise defined.
    \item The transpose of a matrix is \(\top\), e.g.,
        \(\vA^{\top}\). The adjoint is \(\adj\), e.g., \(\vA\adj\).
        The pseudo-inverse is \(\dagger\), e.g., \(\vA^{\dagger}\).
    \item The rank of a matrix is \(\rank(\vA)\), the trace is
        \(\tr(\vA)\), the determinant is \(\det(\vA)\), and the
        log-determinant is \(\logdet(\vA)\).
    \item The element-wise multiplication of matrices is \(\vA \hada
        \vB\). Element-wise squaring is
        \(\vA^{\hada 2}\), etc. Similarly, the Kronecker product of
        matrices is \(\vA \kron \vB\). The iterated Kronecker product
        is \(\vA^{\kron 2}\), etc.
    \item For a function \(f \colon \R \to \R\), its element-wise
        application to a matrix \(\vX\) is \(f[\vX]\), i.e., with
        square brackets. For a function \(f \colon \R^{d} \to
        \R^{k}\) and a matrix \(\vX \in \R^{d \times n}\), we may
        broadcast \(f\) to apply to each column without special
        notation, i.e., \(f(\vX) = [f(\vx_{1}), \dots, f(\vx_{n})]
        \in \R^{k \times n}\).
    \item The (Euclidean) unit sphere is \(\Sphere^{n - 1} \subseteq
        \R^{n}\). The (Euclidean) unit ball is \(\Ball^{n} \subseteq \R^{n}\).
    \item The set of orthogonal matrices are \(\O(m, n) \subseteq
        \R^{m \times n}\) or \(\O(n) = \O(n, n)\).
    \item Symmetric matrices are \(\Sym(n)\), symmetric PSD matrices
        are \(\PSD(n)\), symmetric PD matrices are \(\PD(n)\).
    \item The eigenvalues of symmetric matrix \(\vS \in \Sym(n)\) are
        all real by the spectral theorem; they are denoted
        \(\lambda_{i}(\vS)\) and ordered such that \(\lambda_{1}(\vS)
        \geq \cdots \geq \lambda_{n}(\vS)\).
    \item The singular values of \(\vA \in \R^{m \times n}\) with
        \(\rank(\vA) = r\) are denoted  \(\sigma_{i}(\vA)\) and
        ordered such that \(\sigma_{1}(\vA) \geq \cdots \geq
        \sigma_{r}(\vA) > 0\). By convention we set \(\sigma_{r +
        1}(\vA) = \sigma_{r + 2}(\vA) = \cdots = 0\).
    \item The projection onto a set \(\cK\) is denoted \(\proj_{\cK}\).
    \item The \(\ell^{p}\) norms of vectors are denoted by
        \(\norm{\vx}_{p}\), \(p \in [0, \infty]\).
    \item The Euclidean operator norm on matrices is \(\norm{\vA} =
        \sigma_{1}(\vA)\).
    \item The Frobenius norm on matrices is \(\norm{\vA}_{F} =
        \tr(\vA^{\top}\vA)\).
    \item The Euclidean inner product of vectors is \(\ip{\vx}{\vy} =
        \vy^{\top} \vx\).
    \item The Frobenius inner product on matrices is \(\ip{\vX}{\vY}
        = \tr(\vY^{\top}\vX)\).
    \item The all-ones vector/matrix is denoted \(\vone\) (the shape
        should be obvious from context). Similarly, the all-zeros
        vector/matrix is denoted \(\vzero\). The identity matrix is
        denoted \(\vI\).
\end{itemize}

\section*{Probability Notation}

\begin{itemize}
    \item The base probability measure is \(\Pr\), i.e., the
        probability of an event \(A\) occurring is \(\Pr[A]\). We may
        specify the distribution of a random variable using a
        subscript, i.e., \(\Pr_{\vx \sim \mu}[\vx \in S]\) describes the
        probabilities associated to a random variable $\vx$ with distribution
        (measure) $\mu$. If two random variables $\vx$ and $\vx'$ have the same
        distribution, we write $\vx \equid \vx'$.
    \item The expected value operator is \(\Ex\), with the same
        subscript caveat.
    \item The covariance operator is \(\Cov\), with the same subscript caveat.
    \item The correlation operator is \(\Corr\), with the same subscript caveat.
    \item In certain probabilistic modeling contexts where
        a density is assumed (especially in
        \Cref{ch:denoising,ch:conditional-inference}), we also use the notation
        $p_{\vx}$ for the density of a random variable $\vx$. This extends to
        conditional densities, e.g. $p_{\vx \mid \vy}$ for the density of $\vx$
        conditioned on $\vy$.
    \item We generally use Greek letters to denote
        realizations of random variables (in contrast to the
        random variables themselves). For example, $p_{\vx}(\vxi)$, $p_{\vx \mid
        \vy}(\vxi \mid \vnu)$, $\Ex[\vx \mid \vy=\vnu]$, etc.
    \item The set of probability distributions on a finite set %
        %\footnote{This
        %    should be a measurable space to avoid measure-theoretic and
        %    set-theoretic issues. If you have not studied measure theory,
        %you may relatively safely ignore this.} 
        \(\cX\) is denoted
        \(\Delta(\cX)\). (For \(\cX = [n]\), this is the probability
        simplex which can be embedded into \(\R^{n}\)).
    \item The Gaussian distribution with mean \(\vmu\) and covariance
        \(\vSigma\) is denoted \(\dNorm(\vmu, \vSigma)\).
    \item The uniform distribution over a compact set \(\cX\) is
        denoted \(\dUnif(\cX)\).
\end{itemize}



\section*{Machine Learning Notation}

\begin{itemize}
    \item Encoders and denoisers are usually denoted by \(f\) and
        usually have parameters \(\theta \in \Theta\). We usually
        write the features \(\vz = f_{\theta}(\vx)\).
    \item Decoders are usually denoted by \(g\) and usually have
        parameters \(\eta\). We usually write the auto-encoding
        \(\hat{\vx} = g_{\eta}(\vz)\) corresponding to \(\vx\).
    \item When \(f\) and \(g\) are implemented as neural networks,
        they will have the same number of layers without loss of
        generality; we write them as \(f = f^{L} \circ f^{L - 1}
        \circ \cdots \circ f^{2} \circ f^{1} \circ f^{\pre}\) and \(g
            = g^{\post} \circ g^{L} \circ g^{ L - 1} \circ \cdots \circ
        g^{2} \circ g^{1}\).
    \item We write the features at the input to layer \(\ell\) as
        \(\vz^{\ell}\) such that \(\vz^{\ell + 1} =
        f^{\ell}(\vz^{\ell})\) with \(\vz^{1} = f^{\pre}(\vx)\) and
        \(\vz = \vz^{L + 1}\). Similarly the autoencoding features
        are \(\hat{\vx}^{\ell}\) with \(\hat{\vx}^{\ell + 1} =
        g^{\ell}(\hat{\vx}^{\ell})\) with \(\hat{\vx}^{1} = \vz\) and
        \(\hat{\vx} = g^{\post}(\hat{\vx}^{L + 1})\).
    \item For denoising, optimization, and processes which have a
        continuous time index, the time is almost always a subscript,
        e.g., \(\vx_{t}\). Discrete sequences can have the index be a
        superscript or subscript.
\end{itemize}

\section*{Modeling and Optimization Notation}
\begin{itemize}
    \item As in the Probability and Machine Learning Notation sections above,
        when we have a model, say for the realizations of
        a data distribution supported on $\bR^D$, we always denote it with
        ``plain'' accented variables, say $\vx$.
    \item In cases where we assume our data $\vx$ is generated by a model from
        a restricted class of \textit{parametric models}, say with
        a parameter $\vmu$, we also denote the \textit{true parameter} with
        plain accenting.\footnote{This could arise either due to us making an
        assumption about our data, or due to our data being synthetically
        generated. This convention allows us to adopt a concise notation across
        both `probabilistic' and `analytical' modeling setups. Readers more
        familiar with the latter kinds of frameworks may be used to using
        notation such as $\vmu_{\sharp}$ or $\vmu_o$ for this case; contrast
        with our accented decision variable notation.}
    \item \textit{Decision variables} (or learnable parameters) of a model fit
        to observed data are denoted with ``tilde'' accenting, e.g.\
        $\tilde{\vx}$, in cases where there is an underlying model $\vx$, or
        with plain notation in cases where no confusion is possible (e.g., if
        an empirically-fit model for \textit{nonparametric} data $\vx$ involves
        ``mean'' parameters $\vmu$).
    \item \textit{Optimal solutions} to optimization problems are denoted with
        ``star'' accenting, either as a superscript or a subscript (say
        $\vx_\star$ or $\vx^\star$).
    \item \textit{Estimators} for data that correspond to a specific statistical
        model, especially for the mean of that statistical model, are denoted
        with ``bar'' accenting, say $\bar{\vx}$ (especially for minimum
        mean-squared error denoising in \Cref{ch:denoising}).
    \item \textit{Approximations} associated to computational procedures for
        modeling data are denoted with ``hat'' accenting, say $\hat{\vx}$ for
        the output of an autoencoder or a generative model that approximates
        data $\vx$.\footnote{We generally distinguish these cases from optimal
        solutions $\vx_{\star}$, although in some special cases they are
        equal (e.g., in several examples in \Cref{ch:classic-models}). Think of
        the distinction between the parameter $\vmu$ of a model, which might be
        fit with optimization, and the results of sampling with that learned
        parameter.}
\end{itemize}
